\documentclass{beamer}

% Alternative to \usetheme{claw}
\usepackage{template/beamercolorthemeclaw}
\usepackage{template/beamerfontthemeclaw}
\usepackage{template/beamerinnerthemeclaw}
\usepackage{template/beamerouterthemeclaw}

\usepackage{amsmath, amssymb, amsfonts}% mathematical symbols and the like
\usepackage{amsthm}% definitions, theorems, etc.
\usepackage[colorinlistoftodos]{todonotes}% marking open todos in text/on margins
\usepackage{subfig}% multi-part figures with separate captions per part
\usepackage{url}% render URLs correctly and make them clickable through the hyperref package
\usepackage{longtable}% tables that span multiple pages
\usepackage{booktabs}% tables that actually look good
\usepackage[nolist]{acronym}% consistently use acronyms

\usepackage{listings}
%\usepackage[clean]{svg}



% 1st used for slides footer, 2nd for title slide & PDF metadata
\title[Basilisk]{Basilisk -- Continuous Benchmarking for Triplestores}
% Used for title slide & PDF metadata
\subtitle{}

\date{\today}
% Date with conference name
%\date[\today]{Conference on Fabulous Presentations, 2021}

% One author
% 1st used for slides footer, 2nd for title slide & PDF metadata
\author[Fabian~Rensing]{Fabian Rensing}

% 1st used for slides footer, 2nd for title slide
%\institute[DICE]{DICE Group\\Paderborn University}
% Institute with logo
\institute[]{Supervisor: Prof. Dr. Axel Ngonga\\Paderborn University}

% Two authors from different institutes
% Source: https://www.overleaf.com/learn/latex/Beamer#The_title_page
%\author[\underline{Wilke}, Ngonga]{\underline{Adrian~Wilke}\inst{1} \and Axel~Ngonga\inst{2}}
%\institute[]
%{
%	\inst{1}%
%	DICE Group\\
%	Department of Computer Science\\
%	Paderborn University, Germany
%	\and
%	\inst{2}%
%	DICE Group\\
%	Department of Computer Science\\
%	Paderborn University, Germany
%}

% One logo on title page
%\titlegraphic{\includegraphics[height=1cm]{DICE}}

% Two logos on title page
%\titlegraphic{
%	\includegraphics[height=1cm]{UPB}%
%	\hspace*{2cm}~%
%	\includegraphics[height=1cm]{DICE}%
%}



\begin{document}

% Title page without background
%{\usebackgroundtemplate{}
%\frame[plain]{\titlepage}}

% Title page with background
{\usebackgroundtemplate{\includegraphics[width=\paperwidth]{background-title-upb}
% German UPB logo
%\usebackgroundtemplate{\includegraphics[width=\paperwidth]{background-title-upb-german}
}
\frame[plain]{\titlepage}}

% Begin counting at second frame
\addtocounter{framenumber}{-1}



%\section{Motivation \& Contents}
\begin{frame}{Agenda}
%	\begin{columns}[T]
%	\column{0.45\textwidth}
%		\heading{Motivation}
%		After the title slide, typically a table of contents is presented. Make it more interesting by firstly \emph{introducing the problem} you are solving afterwards. That could also be on an own slide. Better than a long text like this is an image or keywords. Almost always.
%	\column{0.55\textwidth}
		\heading{Contents}\vspace*{.2cm}
		\hypersetup{linkcolor=textblue}
		\tableofcontents
%	\end{columns}
\end{frame}


\chapter{Description}
\label{ch:description}


- halbe Seite
\chapter{Motivation}
\label{ch:motivation}

The motivation

The topic I have to offer has the working title is "Basilisk -- Continuous Benchmarking for Triplestores". At the core, it is developing and deploying a CI/CD tool that hooks into github and/or docker registries.

Described in more detail:

Triplestores -- the database backend of knowledge graphs -- are
typically developed in long iterations and are bench-marked -- if at
all -- only in a very late stage of such an iteration. Benchmarking
and evaluation of benchmarking results are typically done manually
and binds developer's time. Thus, performance regressions are found
very late or never.

With Basilisk we started to develop a continuous benchmarking
platform for triplestore which hooks into github and docker image
registries.

On events like pull requests or newly published versions of
triplestores, a benchmarking suite is run automatically.

The first version of Basilisk
(https://github.com/dice-group/Basilisk,
https://github.com/dice-group/basilisk-frontend) is already
implemented. It is based on the benchmarking tool IGUANA
(https://github.com/dice-group/IGUANA) and Docker. (It requires
triple stores to be dockerized).

The thesis task is to:

1. describe and review the software architecture
2. deploy Basilisk and its frontend on a publicly available VM
3. benchmark 2 versions of Tentris
(https://github.com/dice-group/tentris, via a github hook) and
one version from another triple store (via a docker image
registry, e.g. https://hub.docker.com/r/ontotext/graphdb/ or
https://hub.docker.com/r/openlink/virtuoso-opensource-7).
4. fix critical bugs in 1.-3. and document non-critical

The thesis can be extended to a paper in a scientific journal (e.g.
ISWC) on the resource track with you as first author. 
\chapter{Formulation of Target Setting}
\label{ch:target_setting}
\chapter{Tasks Descriptions}
\label{ch:tasks}

1. describe and review the software architecture
2. deploy Basilisk and its frontend on a publicly available VM
3. benchmark 2 versions of Tentris
(https://github.com/dice-group/tentris, via a github hook) and
one version from another triple store (via a docker image
registry, e.g. https://hub.docker.com/r/ontotext/graphdb/ or
https://hub.docker.com/r/openlink/virtuoso-opensource-7).
4. fix critical bugs in 1.-3. and document non-critical
\chapter{Schedule}
\label{ch:schedule}
\chapter{Preliminary Outline of the Thesis}
\label{ch:outline}

The thesis will comprise the following chapters
\begin{enumerate}
	\item Introduction:
		\begin{itemize}
			\item Description of the topic and the context in which it is inserted
			\item Topic motivation
		\end{itemize}
	
	\item Related Work
	\begin{itemize}
		\item Description of existing techniques and approaches
	\end{itemize}
	
	\item Background
	\begin{itemize}
		\item Description of foundational knowledge required to understand the thesis
	\end{itemize}
	
	\item Approach
	\begin{itemize}
		\item Description of the designed system
	\end{itemize}
	
	\item Implementation
	\begin{itemize}
		\item Implementations details
	\end{itemize}
	
	\item Evaluation
	\begin{itemize}
		\item Experiment setup, requirements and result evaluation
	\end{itemize}
	
	\item Summary and Discussion
	\begin{itemize}
		\item Summary of the work
		\item Highlighting the key findings of the evaluation stage
	\end{itemize}
\end{enumerate}





%%%%%%%%%%%%%%%%%%%%%%%%%%%%%%%%%
%\begin{frame}
%	
%\end{frame}
%
%%\section{Quick Start}
%\begin{frame}[fragile]{Quick Start}
%	Create your first slide:
%	\begin{enumerate}
%		\item Copy all \emph{*.sty} files into a directory
%		\item Copy \emph{packages.tex} into the directory
%		\item Create a .tex file and add the code listed below
%		\item Generate your slide using LaTeX
%	\end{enumerate}
%	\lstset{language=TeX, escapechar=\@, keywordstyle=\color{templatetext}, numbers=none, emph={claw,packages,tex}, emphstyle=\color{textblue}}
%	\begin{lstlisting}[caption={Minimal Example}]
%\documentclass{beamer}
%\usetheme{claw}
%\usepackage{amsmath, amssymb, amsfonts}% mathematical symbols and the like
\usepackage{amsthm}% definitions, theorems, etc.
\usepackage[colorinlistoftodos]{todonotes}% marking open todos in text/on margins
\usepackage{subfig}% multi-part figures with separate captions per part
\usepackage{url}% render URLs correctly and make them clickable through the hyperref package
\usepackage{longtable}% tables that span multiple pages
\usepackage{booktabs}% tables that actually look good
\usepackage[nolist]{acronym}% consistently use acronyms

\usepackage{listings}
%\usepackage[clean]{svg}
%\begin{document}
%\begin{frame} Hello World @\textbackslash@end{frame}
%@\textbackslash@end{document}
%	\end{lstlisting}
%\end{frame}
%
%%\section{Text Formatting}
%%\subsection{Predefined Styles}
%\begin{frame}{Text Formatting}{Predefined Styles}
%	\begin{itemize}
%		\item You could \emph{emphasize} important parts \\(Maybe distinguish between \bad{problems} and \good{solutions})
%		\item Use alert to display \alert{warnings}
%		\item Use the url command (\url{https://dice-research.org/}) or the href command (\href{https://dice-research.org/}{DICE}) for links
%		\item Highlight ``\term{predefined terms}'' like brands and\\ \tech{TechnicalTerms} like software components
%	\end{itemize}
%\end{frame}
%
%%\subsection{Additional Commands}
%\begin{frame}{Text Formatting}{Additional Commands}
%	Use combinations for other concepts:
%	\begin{itemize}
%		\item Text styles: \textbf{bold}, \textit{italic}, \underline{underlined}, \textsc{small caps}
%		\item Font families: \texttt{monospaced}, \textsf{sans serif}, \textrm{roman}
%		\item Text colors:
%			\bluedark{bluedark},
%			\gray{gray},
%			\magenta{magenta},
%			\blue{blue},
%			\orange{orange},
%			\purple{purple},
%			\red{red},
%			\turquoise{turquoise},
%			\green{green}
%		\item Text sizes: {\tiny tiny}, {\scriptsize scriptsize}, {\footnotesize footnotesize}, {\small small}, {\normalsize normalsize}, {\large large}, {\Large Large}, {\LARGE LARGE}, {\huge huge}, {\Huge Huge}
%	\end{itemize}
%\end{frame}
%
%%\section{Code Listings \& Frame Arguments}
%\begin{frame}[fragile]{Code Listings \& Frame Arguments}
%	Use these arguments to configure frames:
%	\setbeamercolor{local structure}{fg=textpurple}
%	\begin{description}
%		\item[\texttt{fragile}\hfill]
%		Specially interpreted contents, e.g. for listings
%		\item[\texttt{plain}\hfill]
%		No headlines, footlines, sidebars; e.g. for large images\\
%		\gray{To also remove background images use:
%		\texttt{\{\textbackslash usebackgroundtemplate\{\}[...]\}}}
%		\item[\texttt{squeeze}\hfill]
%		Squeezes vertical spaces, e.g. for long contents
%		\item[\texttt{shrink}\hfill]
%		Shrinks frame, e.g. for long contents
%	\end{description}
%	\lstset{language=TeX, escapechar=\@, morekeywords={begin,end,frame}, emph={fragile}}
%	\begin{lstlisting}[caption={Frame Options}]
%\begin{frame}[fragile]{Code Listings \& Frame Arguments}
%	% [...]
%@\textbackslash@end{frame}
%	\end{lstlisting}
%\end{frame}
%
%%\section{Mathematics \& Miscellaneous}
%\begin{frame}{Mathematics \& Miscellaneous}
%	\begin{itemize}
%		\item Math\footnote[frame]{This is a footnote also working in columns}: $5^{2}=3^{2}+4^{2}$
%		\item Equations:
%	\end{itemize}
%	\begin{equation}
%		\sum_{n = 1}^{\infty} \frac{1}{n} = 1 + \frac{1}{2} + \frac{1}{3} + \frac{1}{4} + \frac{1}{5} + \dots 
%	\end{equation}
%\end{frame}
%
%%\section{Blocks}
%\begin{frame}[shrink]{Blocks}
%	\begin{block}{This is a Block}
%		\begin{itemize}
%			\item This is an item
%		\end{itemize}
%		\begin{enumerate}
%			\item This is enumeration item
%		\end{enumerate}
%	\end{block}
%	\begin{exampleblock}{This is an Example Block}
%		\begin{itemize}
%			\item This is an item
%		\end{itemize}
%		\begin{enumerate}
%			\item This is enumeration item
%		\end{enumerate}
%	\end{exampleblock}
%	\begin{alertblock}{This is an Alert Block}
%		\begin{itemize}
%			\item This is an item
%		\end{itemize}
%		\begin{enumerate}
%			\item This is enumeration item
%		\end{enumerate}
%	\end{alertblock}
%\end{frame}
%
%% Examples for alternative colors. Configure this in packages.tex, not here.
%\setbeamercolor{block title}{bg=lightorange,fg=templatetext}
%\setbeamercolor{block body}{bg=backgroundorange}
%\setbeamercolor{block title example}{bg=lightgreen,fg=templatetext}
%\setbeamercolor{block body example}{bg=backgroundgreen}
%\setbeamercolor{block title alerted}{bg=lightred,fg=templatetext}
%\setbeamercolor{block body alerted}{bg=backgroundred}
%\begin{frame}[shrink]{Blocks}
%	\begin{block}{This is a Block}
%		\begin{itemize}
%			\item This is an item
%		\end{itemize}
%		\begin{enumerate}
%			\item This is enumeration item
%		\end{enumerate}
%	\end{block}
%	\begin{exampleblock}{This is an Example Block}
%		\begin{itemize}
%			\item This is an item
%		\end{itemize}
%		\begin{enumerate}
%			\item This is enumeration item
%		\end{enumerate}
%	\end{exampleblock}
%	\begin{alertblock}{This is an Alert Block}
%		\begin{itemize}
%			\item This is an item
%		\end{itemize}
%		\begin{enumerate}
%			\item This is enumeration item
%		\end{enumerate}
%	\end{alertblock}
%\end{frame}
%
%%\section{Tables}
%\begin{frame}{Tables}
%	\begin{table}[]
%	\centering
%	\renewcommand{\arraystretch}{1.1}
%	\begin{tabular}{ll}
%	\hline
%	\textbf{Topic} & \textbf{Content} \\ \hline
%	Generator & Use tools like \href{https://www.tablesgenerator.com/}{tablesgenerator.com} \\
%	Large tables & Try the frame option \texttt{[shrink=.8]} \\
%	 & (center table with \texttt{{\textbackslash}hspace*\{5cm\}})\\
%	Large tables & Combine the \href{https://ctan.org/pkg/longtable}{longtable} package and \\
%	 & the frame option \texttt{[allowframebreaks]}\\
%	Style & Try the \href{https://ctan.org/pkg/booktabs}{booktabs} package \\ \hline
%	\end{tabular}
%	\end{table}
%\end{frame}
%
%%\section{Graphs}
%\begin{frame}{Graphs}
%	\centering
%	\begin{tikzpicture}[node distance=1cm and 2cm] 
%		\node[resource] (pb) {Subject}; 
%		\node[literal, right=of pb] (code) {Object};
%		\draw[->] (pb) -- node[above]{Predicate} (code);
%	\end{tikzpicture}
%\end{frame}
%
%\begin{frame}{Thank you!}
%	\centering
%	\vspace{2.5cm}
%	{\huge\textbf{Questions?}}
%	\\\vspace{2.5cm}
%	{\large Data Science Group at Paderborn University}
%	\\\vspace{.1cm}
%	\centering
%	\begin{tabular}{rl}
%	Web: & \href{https://dice-research.org/}{dice-research.org} \\
%	Code: & \href{https://github.com/dice-group}{github.com/dice-group} \\
%	Twitter: & \href{https://twitter.com/DiceResearch}{@DiceResearch}
%	\end{tabular}
%\end{frame}
%
%\section{Appendix \& References}
%\appendix 
%
%\begin{frame}[t,allowframebreaks]{References}
%	\scriptsize
%	\nocite{*}
%	\bibliographystyle{ieeetr}
%	\bibliography{bibliography}
%\end{frame}
%
%\begin{frame}{Appendix}{Predefined Base Colors}
%	\begin{itemize}
%		\item {\color{primarybluedark}\rule{.7cm}{.4cm} primarybluedark}
%		\item {\color{primarybluelight}\rule{.7cm}{.4cm} primarybluelight}
%		\item {\color{primarygraylight}\rule{.7cm}{.4cm} primarygraylight}
%		\item {\color{primarygraydark}\rule{.7cm}{.4cm} primarygraydark}
%		\item {\color{secondarymagenta}\rule{.7cm}{.4cm} secondarymagenta}
%		\item {\color{secondaryblue}\rule{.7cm}{.4cm} secondaryblue}
%		\item {\color{secondarygreen}\rule{.7cm}{.4cm} secondarygreen}
%		\item {\color{secondaryorange}\rule{.7cm}{.4cm} secondaryorange}
%		\item {\color{secondarypurple}\rule{.7cm}{.4cm} secondarypurple}
%		\item {\color{activeyellow}\rule{.7cm}{.4cm} activeyellow}
%		\item {\color{activered}\rule{.7cm}{.4cm} activered}
%		\item {\color{activeturquoise}\rule{.7cm}{.4cm} activeturquoise}
%		\item {\color{activegreen}\rule{.7cm}{.4cm} activegreen}
%		\item {\color{specificblue}\rule{.7cm}{.4cm} specificblue}
%	\end{itemize}
%\end{frame}
%
%\begin{frame}{Appendix}{Predefined Text Colors}
%	\begin{itemize}
%		\item {\color{textbluedark}\rule{.7cm}{.4cm} textdarkblue}
%		\item {\color{textgray}\rule{.7cm}{.4cm} textgray}
%		\item {\color{textmagenta}\rule{.7cm}{.4cm} textmagenta}
%		\item {\color{textblue}\rule{.7cm}{.4cm} textblue}
%		\item {\color{textorange}\rule{.7cm}{.4cm} textorange}
%		\item {\color{textpurple}\rule{.7cm}{.4cm} textpurple}
%		\item {\color{textred}\rule{.7cm}{.4cm} textred}
%		\item {\color{textturquoise}\rule{.7cm}{.4cm} textturquoise}
%		\item {\color{textgreen}\rule{.7cm}{.4cm} textgreen}
%		\item {\color{textbluespecific}\rule{.7cm}{.4cm} textbluespecific}
%	\end{itemize}
%\end{frame}
%
%\begin{frame}{Appendix}{Predefined Element Colors}
%	\begin{itemize}
%		\item {\color{elementgray}\rule{.7cm}{.4cm} elementgray}
%		\item {\color{elementmagenta}\rule{.7cm}{.4cm} elementmagenta}
%		\item {\color{elementblue}\rule{.7cm}{.4cm} elementblue}
%		\item {\color{elementorange}\rule{.7cm}{.4cm} elementorange}
%		\item {\color{elementpurple}\rule{.7cm}{.4cm} elementpurple}
%		\item {\color{elementyellow}\rule{.7cm}{.4cm} elementyellow}
%		\item {\color{elementred}\rule{.7cm}{.4cm} elementred}
%		\item {\color{elementturquoise}\rule{.7cm}{.4cm} elementturquoise}
%		\item {\color{elementgreen}\rule{.7cm}{.4cm} elementgreen}
%		\item {\color{elementbluespecific}\rule{.7cm}{.4cm} elementbluespecific}
%	\end{itemize}
%\end{frame}
%
%\begin{frame}{Appendix}{Predefined Light Colors}
%	\begin{itemize}
%		\item {\color{lightgray}\rule{.7cm}{.4cm} lightgray}
%		\item {\color{lightmagenta}\rule{.7cm}{.4cm} lightmagenta}
%		\item {\color{lightblue}\rule{.7cm}{.4cm} lightblue}
%		\item {\color{lightorange}\rule{.7cm}{.4cm} lightorange}
%		\item {\color{lightpurple}\rule{.7cm}{.4cm} lightpurple}
%		\item {\color{lightyellow}\rule{.7cm}{.4cm} lightyellow}
%		\item {\color{lightred}\rule{.7cm}{.4cm} lightred}
%		\item {\color{lightturquoise}\rule{.7cm}{.4cm} lightturquoise}
%		\item {\color{lightgreen}\rule{.7cm}{.4cm} lightgreen}
%		\item {\color{lightbluespecific}\rule{.7cm}{.4cm} lightbluespecific}
%	\end{itemize}
%\end{frame}
%
%\begin{frame}{Appendix}{Predefined Background Colors}
%	\begin{itemize}
%		\item {\color{backgroundgray}\rule{.7cm}{.4cm}} backgroundgray
%		\item {\color{backgroundmagenta}\rule{.7cm}{.4cm}} backgroundmagenta
%		\item {\color{backgroundblue}\rule{.7cm}{.4cm}} backgroundblue
%		\item {\color{backgroundorange}\rule{.7cm}{.4cm}} backgroundorange
%		\item {\color{backgroundpurple}\rule{.7cm}{.4cm}} backgroundpurple
%		\item {\color{backgroundyellow}\rule{.7cm}{.4cm}} backgroundyellow
%		\item {\color{backgroundred}\rule{.7cm}{.4cm}} backgroundred
%		\item {\color{backgroundturquoise}\rule{.7cm}{.4cm}} backgroundturquoise
%		\item {\color{backgroundgreen}\rule{.7cm}{.4cm}} backgroundgreen
%		\item {\color{backgroundbluespecific}\rule{.7cm}{.4cm}} backgroundbluespecific
%	\end{itemize}
%\end{frame}

\end{document}